Put \(K=\mathcal K_D(m_D)\) and \(g=g_{\rho,\sigma,D}\).  The normalization row of \(A_D\) gives \(\mathbf 1^\top q=1\) for every \(q\in K\).  Thus the assumed nonempty set \(K\) is a closed subset of the finite probability simplex, hence compact, and its linear minimum and maximum are attained.

Apply \(\mathrm{lem:finite\text{-}linear\text{-}programming\text{-}strong\text{-}duality}\) with primal variable \(q\), equality constraint \(A_Dq=m_D\), nonnegativity constraint \(q\geq0\), and cost row \(g\).  The equality multiplier \(\eta\) is unrestricted, while the nonnegative primal coordinates give the dual constraint \(A_D^\top\eta\leq g^\top\).  Therefore
\[
\Gamma_{\rho\prec\sigma}(m_D)
=\min_{q\in K}gq
=\max_{\eta:A_D^\top\eta\leq g^\top}m_D^\top\eta.
\]
For every primal-feasible \(q\) and dual-feasible \(\eta\), weak duality is also the direct calculation
\[
m_D^\top\eta
=q^\top A_D^\top\eta
\leq gq.
\tag{2}
\]

Let \(\mathbf M\in\mathcal C(\mathcal P)\).  The broad projection theorem gives \(q=q_D^\star(\mathbf M)\in K\).  If a feasible \(\eta\) has \(m_D^\top\eta>0\), (2) gives \(gq>0\).  The stated denominator condition gives \(d_Dq>0\), so \(\mathrm{lem:district\text{-}cell\text{-}linearity}\) and the definition \(g=n_{\sigma,D}-n_{\rho,D}\) give
\[
Q_\sigma^\star(\mathbf M)-Q_\rho^\star(\mathbf M)
=\frac{n_{\sigma,D}q-n_{\rho,D}q}{d_Dq}
=\frac{gq}{d_Dq}>0.
\]
Thus every positive dual certificate is sound on the broad fiber.  Division occurs only after the assumed strict positivity has been invoked.

Now also assume \(\mathcal C_{\mathrm{sat},D}(\mathcal P)\ne\varnothing\) and \(\mathsf{ASep}_D(\mathcal P)\).  The focal projection theorem gives \(\mathcal Q_D^{\mathrm{sat},D}(\mathcal P)=K\).  Since the denominator is positive throughout \(K\), uniform strict dominance on the focal fiber is equivalent to \(gq>0\) for every \(q\in K\).  Compactness of \(K\) and continuity of \(q\mapsto gq\) make this equivalent to
\[
\min_{q\in K}gq>0.
\]
Dual attainment makes the last inequality equivalent to existence of a feasible \(\eta\) with \(m_D^\top\eta>0\), proving necessity.

If \(\Gamma_{\rho\prec\sigma}(m_D)<0\), an attained minimizer \(q_-\in K\) has \(gq_-<0\).  Focal projection realizes \(q_-\) by an original-alphabet family in \(\mathcal C_{\mathrm{sat},D}(\mathcal P)\), and positivity of \(d_Dq_-\) gives \(Q_\sigma^\star<Q_\rho^\star\).  If also \(\overline\Gamma_{\rho\prec\sigma}(m_D)>0\), an attained maximizer \(q_+\in K\) is realized in the same way and satisfies \(Q_\sigma^\star>Q_\rho^\star\).  Both families induce \(\mathcal P\); the focal realization changes only district \(D\), so their other districts retain the baseline graph-compatible restrictions.  Structural separation is sufficient here by \(\mathrm{lem:structural\text{-}separation\text{-}implies\text{-}active\text{-}slice}\).  Without focal saturation and active-slice separation, an optimizer of the outer row face need not be induced by a graph-compatible SCM, so no converse lifting is asserted.